\documentclass[b5j,fleqn]{jarticle}
\usepackage{zennyakuUniv}
\begin{document}
\body{
\Title{newcrown3 p59}
\para

\sentP{Jing}{Many Japanese manga and anime are translated into other languages.}{ジン}{多くの日本の漫画やアニメが、ほかの言語に翻訳されています。}

\sentP{Jing}{Sometimes just translating is not enough.}{ジン}{ただ翻訳するだけでは十分でないこともあります。}

\sentP{Hana}{You're right.}{ハナ}{そのとおりですね。}

\sentP{Hana}{There are people who don't know Japanese culture well.}{ハナ}{日本の文化をよく知らない人たちもいます。}

\sentP{Hana}{How do the translators deal with this?}{ハナ}{翻訳者はこのことにどのように対処するのですか。}

\sentP{Jing}{Well, sometimes they explain the culture.}{ジン}{そうですね、文化について説明することもあります。}

\sentP{Jing}{Here's an example from my favorite manga, Nodame Cantabile.}{ジン}{私の大好きな漫画『のだめカンタービレ』から例を一つ紹介します。}

\sentP{Jing}{Look at this page.}{ジン}{このページを見てください。}

\sentP{Hana}{OK.}{ハナ}{わかりました。}

\sentP{Hana}{A woman is eating a bento.}{ハナ}{女性がお弁当を食べています。}

\sentP{Jing}{Right, but readers overseas are unfamiliar with bento.}{ジン}{そうです。しかし、海外の読者にはお弁当になじみがありません。}

\sentP{Jing}{So the translator added an explanation in the back of the book.}{ジン}{そこで翻訳者は、本の巻末に説明を加えました。}
}
\end{document}
